\section{What International AI Collaboration Can and Cannot Fix}
\label{sec:10.3}
AI collaboration reaches furthest where the underlying problem is inherently cross-border and no single country's data is sufficient on its own: climate, where a model trained on one region predicts poorly everywhere else, and disease, where a pathogen does not stay inside the jurisdiction that first detects it. Both have machinery real enough to inspect, and inspecting it shows something other than what the general claim promises.

The institutional machinery in climate is worth describing precisely, because descriptions of it are usually inflated. The Climate, Land, Energy and Water Systems framework — coordinated by UN bodies together with the International Atomic Energy Agency, with a capacity-building program that has reached participants in more than twenty countries — is an integrated systems-optimization framework, and the research literature around it treats machine learning as a promising addition still being built in, not a feature the framework already has \autocite{alexander2025clews}. Where AI collaboration does the specific climate work this section is about, it currently looks more like national weather services training and benchmarking forecasting models against shared data than like a single named research consortium. The named institution is not where the work happens.

Disease surveillance shows the other half. The Global Virome Project, proposed as a successor to USAID's PREDICT program, set out to catalog the viral families with the highest zoonotic spillover potential across dozens of countries before the next pandemic instead of after it \autocite{carroll2018virome}. The AI sits downstream of the cataloging: sequence data of the kind such a project generates is what spillover-risk prediction models get trained on, which makes the international data-sharing arrangement itself, and not any single algorithm, the part actually worth calling collaboration.

Where there is no such arrangement there is nothing to inspect, and the claims get easier rather than harder. Poverty and economic development are where this literature is most confident — analysis revealing patterns of inequality that guides where aid goes, agricultural forecasting, financial services reaching places that have none — and where a specific, verifiable international partnership doing sustained work at that scale is hardest to find.

Neither of the two cases is evidence that AI collaboration works the way the treaty layer aspires to. Each ran on the researchers and institutions involved choosing to participate, project by project, with no enforcement mechanism binding the next one to happen and no guarantee that a government hostile to the work could not simply keep its own researchers out of it. Section~\ref{sec:6.4.1} catalogs what a government willing to point AI at its own population does with tools built for exactly these humanitarian purposes elsewhere. A government can fund pandemic-surveillance research of the kind described above and, separately, build the biometric-surveillance systems that section documents, and nothing about participating in the first constrains the second.
